%% kapa: exam=1 round=23 year=2012 period=2 form=A subject=appraisal_law subject_ko=부동산관계법규 seq=23 number=63
%% source: Q-Net 감정평가사 기출문제 (cst00309) artlSeq=5052580 file=raw/1008/5052580/2012년 제23회 감정평가사 제1차시험 기출문제(2교시 A형).pdf
%% conversion: kordoc 4.24.51 -> tools/gen1.py (auto); review=proofread-sonnet; flags=none
\begin{question}{63}
부동산 가격공시 및 감정평가에 관한 법령상 표준지공시지가에 관한 설명으로 옳은 것은?
\begin{choices}
  \item 국토해양부장관은 표준지에 토지의 사용ㆍ수익을 제한하는 권리가 설정되어 있는 때에는 당해 권리를 포함하여 적정가격을 평가하여야 한다.
  \item 국가가 국유토지를 처분하기 위하여 토지의 가격을 산정하는 경우에 유사한 이용가치를 지닌 표준지공시지가를 적용할 수 없다.
  \item 표준지공시지가에 대하여 이의가 있는 자는 서면으로 시장ㆍ군수 또는 구청장을 거쳐 국토해양부장관에게 이의를 신청하여야 한다.
  \item 국토해양부장관은 표준지공시지가에 대한 이의신청을 받은 날로부터 30일 이내에 이의신청을 심사하여, 그 결과를 신청인에게 서면으로 통지하여야 한다.
  \item 표준지의 적정가격 조사ㆍ평가를 의뢰받은 감정평가업자가 표준지조사평가보고서를 제출하고자 하는 때에는 미리 당해 표준지를 관할하는 시장ㆍ군수 또는 구청장의 의견을 들어야 한다.
\end{choices}
\end{question}
\answer{5}
